\section{Introduction}
\label{sec:intro}

Decision trees are among the most widely used models in machine learning and
data mining~\cite{costa2023survey}. Their predictions are easy to inspect and
explain, so they suit applications in which a model must be understood or
audited. They are also the base learners of random forests and gradient
boosting~\cite{chen2016xgboost,ke2017lightgbm}. CART~\cite{breiman1984cart}
and C4.5~\cite{quinlan1993c45} are two of the best-known algorithms for
inducing classification trees.

Growing a single tree is usually considered cheap. In practice, however, a
tree is rarely trained only once. Tuning the pruning
strength, depth limit or minimum leaf size needs one tree per candidate value,
and $k$-fold cross-validation multiplies this by $k$.
At the same time, datasets with tens of millions of samples or thousands of
attributes arise in areas such as particle physics~\cite{baldi2014higgs},
bioinformatics~\cite{triguero2015ecbdl} and malware
detection~\cite{harang2020sorel}. On such data, a straightforward sequential
implementation can need minutes to hours per tree, and parameter tuning
pays that cost repeatedly.

Modern servers combine many cores on several sockets with high-bandwidth
memory, and many are equipped with data-center GPUs. Using this hardware
efficiently for tree induction is challenging, for two reasons. First, the work per node shrinks with depth. Nodes
near the root hold most of the training set and provide enough work for
thousands of GPU threads. The lower levels consist of many small nodes, for
which kernel-launch, synchronization and transfer costs often dominate.
Second, split evaluation does little computation per value read, so it is
often limited by memory bandwidth.

SLIQ~\cite{mehta1996sliq} and SPRINT~\cite{shafer1996sprint} introduced
presorting to scale to large data. Later work targeted multicore processors~\cite{aldinucci2014fastflow}
and distributed systems~\cite{abuzaid2016yggdrasil}. Several studies proposed
GPU implementations~\cite{lo2014cudt,nasridinov2014gpu,strnad2016gpu}, but
they were evaluated on GPUs of earlier generations. Most current
high-performance GPU tree
learners~\cite{chen2016xgboost,ke2017lightgbm} use
histogram-based split finding: continuous values are assigned to a fixed
number of bins before induction. This is cheaper, but the chosen thresholds, and therefore the trees, usually
differ from those an exact CART or C4.5 search would find. To our knowledge, no
public implementation grows an exact CART or C4.5 tree using the CPU and the
GPU together, with each node sent to one device according to its size.

In this paper, we investigate parallel implementations of CART and C4.5 for
multicore CPUs, CUDA-enabled GPUs and heterogeneous CPU--GPU systems. They
search for splits exactly, without binning: every cut between two consecutive
distinct attribute values is tested. All our implementations grow the same
tree. We checked this node by node, against each other and against
scikit-learn (CART) and Quinlan's original C4.5 program. Our speed comparisons
therefore always compare runs that build the same tree.

We evaluate five implementations, from a straightforward sequential baseline
to a hybrid one that uses the CPU and the GPU together. The experiments use five public datasets with up to
32 million samples and 2,381 attributes, recent CPU servers, and two
generations of data-center GPUs.
% Numbers: Tables tab:time-cart and tab:time-c45 (depth <= 10, no pruning;
% epyc / H100) and bench/results/hybrid-vs-gpu-only (A100/H100, HIGGS, no depth
% limit). Re-derive after the 10-run protocol lands; treat competitor cells
% still marked \est{} as provisional.
The optimized sequential implementation is up to 25$\times$ faster than the
straightforward one. On 48 cores of a two-socket server, the multicore code is
up to 40$\times$ faster again, and the GPU up to 4$\times$ faster than that.
Overall, the best version is up to 1,570$\times$ faster than the
straightforward sequential code. Our multicore implementation trains CART trees up to 360$\times$
faster than scikit-learn~\cite{pedregosa2011sklearn} and C4.5 trees up to
3,100$\times$ faster than Weka~\cite{hall2009weka}. On the GPU, these factors
reach 1,400$\times$ and 8,100$\times$. For deep trees, the hybrid CPU--GPU
implementation is up to 4.9$\times$ faster than growing all nodes on the GPU.

The main contributions of this paper are:
\begin{itemize}
\item exact parallel implementations of CART and C4.5 for multicore CPUs, GPUs
  and heterogeneous CPU--GPU systems, with identical trees across
  implementations;
\item a hybrid CPU--GPU scheme that assigns the nodes of a tree to devices by
  size, so that the CPU and the GPU grow the tree concurrently;
\item an experimental evaluation on large datasets and current hardware,
  covering training time and its breakdown, memory usage, scalability and
  energy consumption, and a comparison with established CPU and GPU
  implementations;
\item an open-source release of the implementations.\footnote{\plan{Repository URL.}}
\end{itemize}

The rest of the paper is organized as follows. Section~\ref{sec:related}
reviews related work. Section~\ref{sec:impl} describes our implementations.
Section~\ref{sec:design} presents the experimental setup and
Section~\ref{sec:results} the results. Section~\ref{sec:discussion} discusses
the findings, and Section~\ref{sec:conclusion} concludes the paper.
